%%%%%%%%%%%%%%%%%%%%%%%%%%%%%%%%%%%%%%%%%%%%%%%%%%%%%%%%%%%%%%%%%
\chapter{SONUÇ VE GELECEK ÇALIŞMALAR}\label{CH6}
%%%%%%%%%%%%%%%%%%%%%%%%%%%%%%%%%%%%%%%%%%%%%%%%%%%%%%%%%%%%%%%%%

\section{Sonuç}\label{sec:sonuc}

Bu tezde, tıbbi görüntünün yalnızca ilgi bölgesini homomorfik olarak şifreleyen seçici şifreleme yaklaşımının izin verilen sızıntısı gerçek tıbbi görüntüler üzerinde ilk kez sistematik olarak değerlendirilmiş ve bu sızıntıyı gideren bir alternatif önerilmiştir. $\Pi_{\mathrm{ROI}}$ protokolü aynı kütüphane ve parametrelerle yeniden üretilmiş, doğruluğu gerçek ROI maskeleriyle doğrulanmış ve maliyeti dört görüntü boyutunda ölçülmüştür.

İki sızıntı istismarı saldırısı, göğüs röntgeni ve beyin MR verisinde şifreli bölgenin klinik içeriğinin büyük ölçüde ele verildiğini göstermiştir. Yalnızca ROI meta verisi teşhisi AUC $0.83$--$0.90$ ile, ROI dışındaki açık bağlam ise AUC $0.97$--$0.99$ ile ele vermiştir. Sızıntı başka bir kaynaktan veriyle eğitilen saldırganda da sürmüş; basit önişleme adımlarıyla ve bağlamın gürültü ya da bulanıklıkla bozulmasıyla giderilememiştir. Adaptif saldırgana karşı değerlendirilen savunmalar, teşhisi gizlemek için görüntünün neredeyse tamamının şifrelenmesi gerektiğini ve bu durumda seçici şifrelemenin hız kazancının tam şifrelemeye göre yaklaşık 1 kata indiğini göstermiştir.

Bu bulgulardan hareketle önerilen Odaklı Tam Şifreleme (FoveaHE), görüntüden ROI merkezli, sabit boyutlu ve çok çözünürlüklü bir temsil çıkarıp temsilin tamamını şifrelemektedir. Deneylerden önce belirlenen beş başarı ölçütünün tamamı 5 tohumla karşılanmıştır. Sunucunun gördüğü bilgilerle kurulan saldırgan şans düzeyinde kalmış (AUC $0.45$--$0.53$), tam şifrelemeye göre beyin MR görüntüsünde 29--37 kat, göğüs röntgeninde 7--9 kat hızlanma elde edilmiş ve teşhis başarısı tam görüntüyle eğitilen aynı model sınıfının üstüne çıkmıştır ($0.881$'e karşı $0.948$ ve $0.946$'ya karşı $0.954$). Aynı ölçütlerle karşılaştırılan beş yaklaşım arasında gizlilik şartlı $\Pi_{\mathrm{ROI}}$ hız kazancını yitirmiş, bozuk bağlam sızdırmaya devam etmiş, önceden eğitilmiş bir ağın özetini şifreleyen yaklaşım ise FoveaHE'den daha doğru bulunmuştur. Bütçe duyarlı odak deneyi, çözünürlüğün odak ile genel bakış arasında paylaştırılmasının her iki uç durumdan da iyi olduğunu; harici derleme testi ise temsiller arasındaki sıralamanın farklı bir derlemede korunduğunu göstermiştir.

Tezin temel sonucu iki cümleyle özetlenebilir: \emph{ROI tabanlı seçici homomorfik şifreleme, tıbbi görüntülerde gizlilik şartı konduğunda verimlilik avantajını kaybetmektedir. Verimlilik açık bırakılan alandan değil, şifrelenen verinin çözünürlüğünden kazanıldığında gizlilik ile verimlilik birlikte sağlanabilmektedir.} Bu sonuç, seçici şifreleme önerilerinin yalnızca verimlilikle değil, izin verdikleri sızıntının gerçek veride ölçülmesiyle birlikte değerlendirilmesi gerektiğini ortaya koymaktadır. Tez bu amaçla yeniden kullanılabilir bir sızıntı denetim hattı, somut tasarım kuralları ve sızıntısız bir alternatif sunmaktadır.

\section{Gelecek Çalışmalar}\label{sec:gelecek}

\begin{itemize}
\item \textbf{Bağımsız veri kümeleri.} Bulgular, hastane bazında ayrılmış ve yetişkin hastalardan oluşan bağımsız göğüs röntgeni kümelerinde tekrarlanmalıdır.
\item \textbf{Odaklı özet.} Odak ve genel bakışın önceden eğitilmiş bir ağla ayrı ayrı özetlenip birlikte şifrelenmesi, FoveaHE ile şifreli özetin güçlü yanlarını birleştirebilir.
\item \textbf{Öğrenilebilir odak.} Odak penceresinin ve çözünürlük paylaşımının görevle birlikte uçtan uca öğrenilmesi ile ROI'nin istemcide hafif bir modelle bulunmasının maliyeti incelenmelidir.
\item \textbf{Daha derin şifreli modeller.} Daha derin şifreli ağlarda ve farklı CKKS parametreleriyle odaklı temsilin doğruluk ve maliyet dengesi ölçülmeli; ``Encrypt What Matters'' çalışmasındaki yerel mimarilerle gizlilik şartlı bir karşılaştırma yapılmalıdır \cite{backour2026encrypt}.
\item \textbf{Üç boyutlu MR verisi.} $\Pi_{\mathrm{ROI}}$ çalışmasının motivasyon kaynağı olan üç boyutlu beyin tümörü MR verisinde hem tümör çevresindeki kitle etkisinin sızıntısı hem de odaklı temsilin üç boyuta genişletilmesi incelenmelidir.
\item \textbf{Hasta yeniden tanımlama.} ROI dışındaki bölgenin hastayı yeniden tanımlamaya yetip yetmediği ölçülmelidir \cite{packhauser2022reid}.
\item \textbf{Federatif öğrenme.} Aynı sızıntı değerlendirme yaklaşımı, federatif öğrenmede gradyan düzeyinde seçici şifrelemeye taşınarak yüksek lisans çalışmasına zemin oluşturulabilir.
\end{itemize}
